\begin{resumo}[RESUMO]
\vspace*{-6mm}

<<Apresentar brevemente o conceito principal do trabalho, como o Cluster API, e destaque o interesse recente da indústria em tecnologias de orquestração e controle distribuído.>> 

<<Explicar os desafios que existem hoje para integrar múltiplos ambientes industriais ou computação de borda, e por que isso exige novas soluções.>> 

Neste contexto, este trabalho propõe uma investigação revisional sobre o uso do Cluster API como mecanismo de controle distribuído em cenários industriais com múltiplos nós de operação, como fábricas remotas ou dispositivos de borda.

<<Descrever a proposta geral do trabalho, indicando que ela envolve tanto um levantamento bibliográfico estruturado quanto uma validação prática em um ambiente experimental.>> 

Para isso, foi construído um ambiente com um nó de controle e múltiplos nós gerenciados em diferentes localizações, permitindo avaliar aspectos como latência, confiabilidade e viabilidade de adoção em ambientes industriais reais. 

<<Resumir as técnicas, ferramentas ou abordagens usadas — ex: tipos de clusters testados, uso de alguma nuvem pública, métricas medidas etc.>> 

Os testes realizados demonstraram <<insira aqui uma síntese dos principais resultados esperados ou obtidos>>. 

Os resultados sugerem que o uso do Cluster API pode representar uma solução promissora para a integração e controle de sistemas industriais distribuídos, especialmente em contextos que exigem escalabilidade, autonomia e baixa latência.

Palavras-chave: Cluster API. Computação de borda. Orquestração de clusters. Indústria 4.0. Integração distribuída.
\end{resumo}
